\begin{center}
{\large\bfseries Evaluation Protocol Outweighs Multimodal Fusion in Wearable Stress and\\[2pt]
Attention Monitoring: Systematic Review and Meta-Analysis}\\[10pt]
\end{center}

\section*{Abstract}
Multimodal fusion of wearable and camera-based sensing is widely presented as the route to
deployable stress and attention monitoring, yet the supporting evidence is assembled from
between-study comparisons that confound fusion with cohort, dataset and evaluation protocol.
We screened 116 records published between 2022 and 2026, retained 60 reporting stress,
attention, drowsiness or cognitive-load outcomes, and independently corroborated every record
before use. Restricting synthesis to within-study contrasts, in which one team holds cohort,
pipeline and protocol fixed, the median fusion gain was 4.67 percentage points (pp; $k=4$,
range 1.00 to 14.30). Across 10 within-study evaluation contrasts the median gap was 9.14 pp,
roughly twice as large: personalization alone reached 27.41 pp of accuracy and 48.66 pp of F1
within one study, and laboratory-to-field transfer cost 13.00 pp of F1. On WESAD, 7 of 9
records exceeded 92\% irrespective of architecture, and leave-one-subject-out results were
indistinguishable from those whose protocol was never stated. Fusion helps, but less than
evaluation design; subject-independent field validation, not further architectural
elaboration, is the binding constraint on deployable systems.

\textbf{Keywords:} stress detection; attention monitoring; multimodal fusion; wearable
biosensors; evaluation protocol.

\section{Introduction}

\subsection{Background and Motivation}
Stress and attention are the two physiological and cognitive states that most directly
govern human performance, safety and long-term health. Acute stress mobilizes the autonomic
nervous system, and sustained or repeated activation degrades sustained attention, working
memory and decision quality. In settings where a lapse is consequential, such as driving,
clinical care, aviation and industrial operation, continuous and objective measurement of
these states has clear practical value. That motivation has driven work in occupational
health \cite{r4,r5,r61}, clinical monitoring \cite{r3,r28,r115}, road safety \cite{r7,r8,r45},
education \cite{r30,r74,r99} and free-living daily life \cite{r10,r13,r35}.

The engineering question, however, has shifted. Enough is now known about which signals carry
stress-related information that the open problem is no longer detection in principle but
detection that survives contact with a new person, a new day and a new environment. This
review is organized around that shift.

\subsection{Sensing Modalities}
Wearable devices sample autonomic activity through electrocardiography (ECG) and the
heart-rate variability (HRV) derived from it, photoplethysmography (PPG) and blood volume
pulse, electrodermal activity (EDA), skin temperature, respiration and
electroencephalography (EEG). Research-grade wristbands \cite{r12,r19}, consumer smartwatches
\cite{r13,r25}, in-ear devices \cite{r43}, behind-the-ear EEG \cite{r31} and custom sensor
packages \cite{r14,r108} now cover most of the practical design space. Inertial measurement
units are almost always co-located, supplying the motion context needed to separate
autonomic arousal from physical exertion \cite{r19,r79}.

Camera-based sensing supplies a complementary, contactless channel. Facial action unit
analysis, gaze and blink dynamics and head pose expose behavioral and affective information
that physiology does not carry, while remote PPG recovers cardiac information from subtle
color changes in ordinary RGB video and so overlaps directly with the wearable channel
\cite{r16}. Systems that combine camera streams with wearable PPG and EDA in a single
pipeline are now reported in workplace settings \cite{r17}. Figure~\ref{fig:block} shows the
canonical processing chain shared by the systems reviewed here, from acquisition through
conditioning, representation, fusion and inference to the decision or closed-loop actuation
stage. The element we place on an equal footing with the others, and which motivates this
review, is the evaluation protocol: it is not a reporting detail appended to a system but the
stage that determines what any reported accuracy actually means.

\begin{figure}[!htbp]
\centering
\resizebox{\textwidth}{!}{\begin{tikzpicture}[
  x=1mm, y=1mm,
  font=\sffamily,
  card/.style={rectangle, rounded corners=1.2mm, draw=#1, fill=#1!6,
               line width=0.55pt, align=center, inner sep=0pt},
  stage/.style={rectangle, rounded corners=1mm, fill=black!5, draw=none,
                inner sep=0pt, minimum width=17mm, align=center,
                font=\sffamily\bfseries\footnotesize, text=black!62},
  chip/.style={rectangle, rounded corners=0.9mm, draw=#1!70, fill=white,
               line width=0.45pt, inner xsep=1.7mm, inner ysep=1.0mm,
               font=\sffamily\footnotesize, text=#1!55!black},
  flow/.style={-{Straight Barb[length=1.5mm, width=1.5mm]},
               line width=0.8pt, draw=black!45},
  loop/.style={-{Straight Barb[length=1.5mm, width=1.5mm]},
               line width=0.7pt, draw=RED!75, dash pattern=on 1.4mm off 1.1mm},
  ttl/.style={font=\sffamily\bfseries\footnotesize, text=black!85},
  sub/.style={font=\sffamily\footnotesize, text=black!58},
]

\definecolor{GRN}{HTML}{2F7A4F}
\definecolor{PUR}{HTML}{7A4FA8}
\definecolor{ORG}{HTML}{C2701C}
\definecolor{BLU}{HTML}{2B6CB0}
\definecolor{RED}{HTML}{C03434}
\definecolor{TEA}{HTML}{2C7A7B}
\definecolor{GLD}{HTML}{A9761A}

%% ---------- geometry ----------
\def\SW{27}   % sensor card width
\def\SH{20}   % sensor card height
\def\GAP{2.6} % gap between sensor cards
\def\LX{-25}  % stage label column x

% mini waveform helper: #1 colour, #2 anchor node, #3 y-offset, #4 coords
\newcommand{\wave}[3]{\draw[line width=0.55pt, draw=#1, line cap=round] #2 #3;}

%% ================= STAGE 1 : ACQUISITION =================
\node[card=GRN, minimum width=\SW mm, minimum height=\SH mm] (s1) at (0,0) {};
\node[card=GRN, minimum width=\SW mm, minimum height=\SH mm] (s2) at (\SW+\GAP,0) {};
\node[card=GRN, minimum width=\SW mm, minimum height=\SH mm] (s3) at (2*\SW+2*\GAP,0) {};
\node[card=GRN, minimum width=\SW mm, minimum height=\SH mm] (s4) at (3*\SW+3*\GAP,0) {};
\node[card=GRN, minimum width=\SW mm, minimum height=\SH mm] (s5) at (4*\SW+4*\GAP,0) {};
\node[card=PUR, minimum width=\SW mm, minimum height=\SH mm] (s6)
      at (5*\SW+5*\GAP+7,0) {};

\foreach \n/\t/\l in {s1/{ECG / HRV}/{chest, patch},
                      s2/{PPG / BVP}/{wrist, in-ear},
                      s3/{EDA, skin temp.}/{wrist, finger},
                      s4/{EEG}/{scalp, behind-ear},
                      s5/{IMU}/{accel., gyro.},
                      s6/{RGB camera}/{face, gaze, pose}} {
  \node[ttl] at ([yshift=6.0mm]\n.center) {\t};
  \node[sub] at ([yshift=2.9mm]\n.center) {\l};
}

% --- signal glyphs ---
\begin{scope}[shift={([yshift=-3.6mm]s1.center)}]   % ECG
  \draw[line width=0.62pt, draw=GRN, line join=round]
    (-12,0) -- (-9.4,0) -- (-8.9,0.55) -- (-8.4,-0.9) -- (-8.05,3.3)
    -- (-7.7,-3.4) -- (-7.3,1.2) -- (-6.7,-0.25) -- (-4.6,0)
    -- (-2.0,0) -- (-1.5,0.55) -- (-1.0,-0.9) -- (-0.65,3.3)
    -- (-0.3,-3.4) -- (0.1,1.2) -- (0.7,-0.25) -- (2.8,0)
    -- (5.4,0) -- (5.9,0.55) -- (6.4,-0.9) -- (6.75,3.3)
    -- (7.1,-3.4) -- (7.5,1.2) -- (8.1,-0.25) -- (12,0);
\end{scope}
\begin{scope}[shift={([yshift=-3.6mm]s2.center)}]   % PPG
  \draw[line width=0.62pt, draw=GRN, smooth]
    plot[domain=-12:12, samples=110]
    (\x, {1.05*sin((\x+12)*45)^2 + 0.42*sin((\x+12)*90) - 0.9});
\end{scope}
\begin{scope}[shift={([yshift=-3.6mm]s3.center)}]   % EDA
  \draw[line width=0.62pt, draw=GRN, smooth]
    plot[domain=-12:12, samples=80]
    (\x, {-2.1 + 0.075*(\x+12) + 2.05*exp(-((\x-1)/3.4)^2)});
\end{scope}
\begin{scope}[shift={([yshift=-3.6mm]s4.center)}]   % EEG
  \draw[line width=0.55pt, draw=GRN, smooth]
    plot[domain=-12:12, samples=180]
    (\x, {0.95*sin(\x*152) + 0.62*sin(\x*63) + 0.34*sin(\x*311)});
\end{scope}
\begin{scope}[shift={([yshift=-3.6mm]s5.center)}]   % IMU
  \draw[line width=0.55pt, draw=GRN]
    (-12,1.3) -- (-8,1.3) -- (-8,2.5) -- (-4,2.5) -- (-4,0.9) -- (0.5,0.9)
    -- (0.5,2.1) -- (5,2.1) -- (5,1.3) -- (12,1.3);
  \draw[line width=0.55pt, draw=GRN!62]
    (-12,-1.0) -- (-6.5,-1.0) -- (-6.5,-2.3) -- (-1.5,-2.3) -- (-1.5,-0.3)
    -- (3.5,-0.3) -- (3.5,-1.6) -- (12,-1.6);
\end{scope}
\begin{scope}[shift={([yshift=-3.7mm]s6.center)}]   % camera: face + gaze cone
  \draw[line width=0.5pt, draw=PUR!55, dash pattern=on 0.7mm off 0.6mm]
        (-5.2,-3.9) rectangle (5.2,4.1);
  \draw[line width=0.62pt, draw=PUR] (-1.6,0) ellipse [x radius=2.9, y radius=3.5];
  \fill[PUR] (-2.65,0.95) circle (0.33);
  \fill[PUR] (-0.55,0.95) circle (0.33);
  \draw[line width=0.45pt, draw=PUR] (-2.7,-1.4) .. controls (-1.6,-2.4) .. (-0.5,-1.4);
  \draw[line width=0.4pt, draw=PUR!55] (-0.55,0.95) -- (4.6,2.5);
  \draw[line width=0.4pt, draw=PUR!55] (-0.55,0.95) -- (4.6,-0.6);
  \fill[PUR!14] (-0.55,0.95) -- (4.6,2.5) -- (4.6,-0.6) -- cycle;
\end{scope}

% group frames
\begin{scope}[on background layer]
  \node[draw=GRN!55, dash pattern=on 1.1mm off 0.9mm, rounded corners=1.5mm,
        line width=0.6pt, fill=GRN!3, fit=(s1)(s5), inner sep=2.4mm] (wgrp) {};
  \node[draw=PUR!55, dash pattern=on 1.1mm off 0.9mm, rounded corners=1.5mm,
        line width=0.6pt, fill=PUR!3, fit=(s6), inner sep=2.4mm] (vgrp) {};
\end{scope}
\node[font=\sffamily\bfseries\footnotesize, text=GRN, above=1.3mm of wgrp.north]
      {CONTACT SENSING: wearable physiological and inertial channels};
\node[font=\sffamily\bfseries\footnotesize, text=PUR, above=1.3mm of vgrp.north]
      {CONTACTLESS};

%% ================= STAGE 2 : CONDITIONING =================
\def\CW{91}
\node[card=ORG, minimum width=\CW mm, minimum height=16mm] (c1) at (32.2,-31) {};
\node[card=ORG, minimum width=78mm, minimum height=16mm] (c2) at (127,-31) {};
\node[ttl] at ([yshift=4.3mm]c1.center) {Signal conditioning};
\node[sub, align=center] at ([yshift=-1.3mm]c1.center)
      {band-pass filtering \ \textbullet\ \ motion-artifact rejection\\
       resampling \ \textbullet\ \ signal-quality gating};
\node[ttl] at ([yshift=4.3mm]c2.center) {Vision front end};
\node[sub, align=center] at ([yshift=-1.3mm]c2.center)
      {face detection \ \textbullet\ \ facial landmarks\\
       action units \ \textbullet\ \ gaze \ \textbullet\ \ rPPG extraction};

%% ================= STAGE 3 : REPRESENTATION =================
\node[card=ORG, minimum width=\CW mm, minimum height=16mm] (r1) at (32.2,-52) {};
\node[card=ORG, minimum width=78mm, minimum height=16mm] (r2) at (127,-52) {};
\node[ttl] at ([yshift=4.3mm]r1.center) {Physiological representation};
\node[sub, align=center] at ([yshift=-1.3mm]r1.center)
      {HRV indices \ \textbullet\ \ EDA phasic and tonic components\\
       time-frequency maps \ \textbullet\ \ learned embeddings};
\node[ttl] at ([yshift=4.3mm]r2.center) {Behavioral representation};
\node[sub, align=center] at ([yshift=-1.3mm]r2.center)
      {action-unit intensities \ \textbullet\ \ blink rate\\
       fixation statistics \ \textbullet\ \ rPPG-derived heart rate};

%% ================= STAGE 4 : FUSION =================
\node[card=TEA, minimum width=179.4mm, minimum height=15mm] (fu) at (76.2,-73) {};
\node[ttl, anchor=west] at ([xshift=4mm,yshift=3.6mm]fu.west) {Fusion stage};
\node[chip=TEA, anchor=west] (f1) at ([xshift=4mm,yshift=-2.6mm]fu.west)
      {feature level (early)};
\node[chip=TEA, right=3mm of f1] (f2) {decision level (late)};
\node[chip=TEA, right=3mm of f2] (f3) {hybrid, multi-stage};
\node[chip=TEA, right=3mm of f3] (f4) {context gated};
\node[sub, anchor=east] at ([xshift=-4mm,yshift=3.6mm]fu.east)
      {median within-study gain \textbf{+4.67 pp}};

%% ================= STAGE 5 : INFERENCE =================
\node[card=BLU, minimum width=179.4mm, minimum height=15mm] (inf) at (76.2,-92) {};
\node[ttl, anchor=west] at ([xshift=4mm,yshift=3.6mm]inf.west) {Inference model};
\node[chip=BLU, anchor=west] (i1) at ([xshift=4mm,yshift=-2.6mm]inf.west)
      {classical ML: RF, XGBoost, SVM};
\node[chip=BLU, right=3mm of i1] (i2) {CNN, LSTM, CNN--LSTM};
\node[chip=BLU, right=3mm of i2] (i3) {transformer, attention};
\node[chip=BLU, right=3mm of i3] (i4) {foundation model};

%% ================= STAGE 6 : EVALUATION (focus) =================
\node[rectangle, rounded corners=1.4mm, draw=RED, fill=RED!7, line width=1.15pt,
      minimum width=179.4mm, minimum height=19mm, inner sep=0pt] (ev) at (76.2,-113) {};
\node[ttl, anchor=west, text=RED!72!black]
      at ([xshift=4mm,yshift=5.2mm]ev.west) {Evaluation protocol};
\node[chip=RED, anchor=west] (e1) at ([xshift=4mm,yshift=-0.6mm]ev.west) {personalized};
\node[chip=RED, right=3mm of e1] (e2) {subject dependent};
\node[chip=RED, right=3mm of e2] (e3) {leave-one-subject-out};
\node[chip=RED, right=3mm of e3] (e4) {cross-dataset};
\node[chip=RED, right=3mm of e4] (e5) {laboratory vs field};
\node[font=\sffamily\itshape\footnotesize, text=RED!72!black, anchor=west]
      at ([xshift=4mm,yshift=-5.9mm]ev.west)
      {this stage, not the fusion stage, sets what a reported accuracy means:
       median within-study swing \textbf{9.14 pp}, up to \textbf{48.66 pp}};

%% ================= STAGE 7 : OUTPUT =================
\node[card=GLD, minimum width=57mm, minimum height=14mm] (o1) at (15,-137) {};
\node[card=GLD, minimum width=57mm, minimum height=14mm] (o2) at (76.2,-137) {};
\node[card=GLD, minimum width=57mm, minimum height=14mm] (o3) at (137.4,-137) {};
\node[ttl] at ([yshift=2.6mm]o1.center) {Stress state and level};
\node[sub] at ([yshift=-1.7mm]o1.center) {binary or graded};
\node[ttl] at ([yshift=2.6mm]o2.center) {Attention and drowsiness};
\node[sub] at ([yshift=-1.7mm]o2.center) {cognitive load, inattention};
\node[ttl] at ([yshift=2.6mm]o3.center) {Closed-loop actuation};
\node[sub] at ([yshift=-1.7mm]o3.center) {feedback, neurostimulation};

%% ================= stage labels =================
\node[stage, minimum height=\SH mm]  at (\LX, 0)    {ACQUIRE};
\node[stage, minimum height=16mm]    at (\LX,-31)   {CONDITION};
\node[stage, minimum height=16mm]    at (\LX,-52)   {REPRESENT};
\node[stage, minimum height=15mm]    at (\LX,-73)   {FUSE};
\node[stage, minimum height=15mm]    at (\LX,-92)   {INFER};
\node[stage, minimum height=19mm, fill=RED!12, text=RED!70!black] at (\LX,-113) {EVALUATE};
\node[stage, minimum height=14mm]    at (\LX,-137)  {DECIDE};

%% ================= flow arrows =================
\draw[flow] (wgrp.south) -- (c1.north);
\draw[flow] (vgrp.south) -- (c2.north);
\draw[flow] (c1) -- (r1);
\draw[flow] (c2) -- (r2);
\draw[flow] (r1.south) -- ++(0,-3) -| ([xshift=-44mm]fu.north);
\draw[flow] (r2.south) -- ++(0,-3) -| ([xshift=44mm]fu.north);
\draw[flow] (fu) -- (inf);
\draw[flow] (inf) -- (ev);
\draw[flow] (ev.south) -- ++(0,-4) -| (o1.north);
\draw[flow] (ev.south) -- ++(0,-4) -| (o2.north);
\draw[flow] (ev.south) -- ++(0,-4) -| (o3.north);
\draw[loop] (o3.east) -- ++(8.5,0) |- ([yshift=-1.5mm]vgrp.east);
\node[font=\sffamily\itshape\footnotesize, text=RED!75, anchor=south, align=center,
      fill=white, inner sep=0.6mm, rotate=90]
      at ([xshift=10mm,yshift=28mm]o3.east) {closed-loop actuation};
\end{tikzpicture}
}
\caption{Seven-stage processing chain shared by the wearable and camera-based stress and
attention monitoring systems reviewed here, with a representative trace drawn inside each
acquisition card. Contact channels (green) and the contactless camera channel (purple) are
conditioned and represented separately before being combined at a feature, decision, hybrid
or context-gated fusion stage and passed to an inference model. The evaluation protocol is
drawn as a stage of equal standing, and highlighted, because it is the stage that determines
what a reported accuracy means: the median within-study swing attributable to it is 9.14
percentage points against 4.67 for the fusion stage. The dashed return path denotes
closed-loop actuation, in which a decision drives stimulation or feedback that the sensors
then observe.}
\label{fig:block}
\end{figure}

\subsection{Limitations of Unimodal Approaches and the Case for Integration}
Each modality fails in characteristic ways. Wearable signals are corrupted by motion
artifacts and variable skin contact, and physiological arousal caused by exertion is not
separable from arousal caused by stress without contextual information \cite{r19,r79}.
Camera-based systems depend on illumination, an unoccluded face and a favorable viewing
angle, and observe only the externally expressed surface of an internal state.

The standard argument is that these failure modes are complementary, so fusing the channels
should be more than additive. The argument is plausible, and several studies do report gains.
What is rarely examined is the magnitude of that gain once it is measured properly. Almost
all quantitative support for fusion in this literature is assembled by comparing the accuracy
reported by a multimodal study against the accuracy reported by a different unimodal study.
Those two numbers differ in cohort, stressor protocol, label source, class definition,
preprocessing and, critically, validation scheme. A between-study difference therefore
estimates the sum of all of these, not the effect of fusion. This review is built on the
observation that a much cleaner estimate is available and has not been systematically
exploited: within a single study, a single team frequently reports a unimodal baseline and a
fused model under an identical protocol, and separately reports the same model under two
different evaluation regimes. These within-study contrasts identify the fusion effect and the
evaluation effect on the same scale, and permit them to be compared directly.

\subsection{Research Questions and Objectives}
Prior reviews of wearable stress monitoring have documented that generalization is weak
\cite{r118}, and prior reviews of affective computing survey physiological and vision-based
modalities largely in parallel. Neither has quantified how the benefit attributed to fusion
compares with the performance swing produced by evaluation design. This review addresses that
gap through three questions:

\textbf{RQ1:} When fusion gain is estimated only from within-study, protocol-matched
comparisons, how large is it? This matters because every deployment decision that favors a
second sensor over a single one is justified by an assumed magnitude, and that magnitude has
never been isolated from its confounds.

\textbf{RQ2:} How does that fusion gain compare with the within-study performance difference
attributable to evaluation choices, namely personalization, deployment setting, dataset and
task granularity? This matters because if the second effect dominates, the field's
optimization target is misplaced.

\textbf{RQ3:} Has the dominant benchmark saturated, and how completely is the evidence base
reported and independently verifiable? This matters because a saturated benchmark combined
with unreported protocols makes architectural progress unfalsifiable.

The review pursues five objectives:
\begin{itemize}[leftmargin=1.25em,itemsep=1pt,topsep=2pt]
\item \textbf{O1:} Assemble a topically screened corpus of 2022 to 2026 studies in which every
      bibliographic record is independently corroborated.
\item \textbf{O2:} Carry only corroborated outcome values into the synthesis and report the
      corroboration and reporting-recovery rates as evidence in their own right.
\item \textbf{O3:} Estimate the within-study fusion gain under matched protocols.
\item \textbf{O4:} Estimate the within-study evaluation-induced gap and compare the two on a
      common scale.
\item \textbf{O5:} Test benchmark saturation and derive a reporting specification for
      translational work.
\end{itemize}

Section~2 reviews unimodal and integrated work and the limits of prior surveys. Section~3
gives the review and verification methodology. Section~4 reports results against the
objectives. Section~5 synthesizes the trends. Sections~6 and~7 discuss and conclude.
